\section{Algorithme universel de la Loi Générale des Masses}
\Needspace{7\baselineskip}
\subsection{Domaine du lecteur cellulaire : extension survivante, parcours, enregistrement et signature}
L'entrée contient uniquement :

\[
 (\mathrm{APP},\mathrm{TRIT},\mathrm{TPK},
 \mathcal R_{324},\mathcal L_{108},
 \text{représentation},\text{couplage},
 \text{carte dimensionnelle commune}).
\]
Elle ne contient ni masse observée, ni résidu, ni coefficient choisi par approximation, ni tableau externe.

\Needspace{7\baselineskip}
\subsection{Construction fonctorielle de la réalisation matérielle : fibre de parcours, entrelaceur, projecteur de coïncidence, opérateur spectral et scalarisation}
Pour chaque état physique \(X\) :

\begin{enumerate}
\item construire son descripteur de charge, génération, saveur, couleur, spin, statistique,
orientation et composition ;

\item former la fibre des parcours compatibles ;
\item calculer l'espace des entrelaceurs
\end{enumerate}
   \[
   \mathscr I_X=
   \operatorname{Hom}_{G_X}(V_X,\ell^2(F_X));
   \]
\begin{enumerate}
\item si \(\dim\mathscr I_X=0\), écarter cette identification ;
\item construire le projecteur par moyenne de Reynolds et facteur polaire ;
\item reconstruire les enregistrements généalogiques de tous les parcours du support ;
\item calculer la signature, \(K_D\), \(K_\Omega\) et
\(\mathcal F_{\rm tower}^{2w}\) ;

\item appliquer le pointeur de lecture du couplage, l'égalisateur ou la mesure
spectrale conjointe ;

\item construire \(\widehat M_X\) ;
\item réduire à un scalaire seulement si la variance est nulle ou s'il existe un pôle isolé ;
\item pour les composés, appliquer d'abord \(\mathfrak C_{12}\) ;
\item pour les largeurs, calculer le semi-groupe de survie ;
\item multiplier par la section dimensionnelle commune ;
\item figer le résultat ;
\item ouvrir la comparaison externe.
\end{enumerate}
\Needspace{7\baselineskip}
\subsection{Existence de la réalisation par fibres}
La multiplicité d'une représentation \(\pi_X\) au sein d'une fibre se calcule par

\[
 \dim\mathscr I_X=
 \frac1{|G_X|}
 \sum_{g\in G_X}
 \overline{\chi_{\pi_X}(g)}
 |\operatorname{Fix}_{F_X}(g)|.
\]
Si elle est positive, il existe un entrelaceur. Si elle est supérieure à un, le résultat est un faisceau de réalisations et le couplage doit sélectionner un état ou conserver la mesure mixte. Cette formule remplace une recherche nominale par une condition algébrique finie sur les 324 parcours.

\Needspace{7\baselineskip}
\subsection{Égalisateur spectral des lecteurs \(K_D\) et \(K_\Omega\) au moyen de \(Q_{\mathrm{eq}}=\mathbf1_{\{0\}}(\widehat L_D-\widehat L_\Omega)\)}
Soient

\[
 L_r(\Gamma)=
 \kappa_r(\Gamma)\log R_{\rm act}
 +\sum_{h\in\mathfrak S}\eta_h^r(\Gamma)s_h,
 \qquad r\in\{D,\Omega\}.
\]
Le projecteur de coïncidence est

\[
 Q_{\rm eq}=
 \mathbf1_{\{0\}}(\widehat L_D-\widehat L_\Omega).
\]
Une réalisation admet une lecture commune exactement lorsque

\[
 P_X\le Q_{\rm eq}.
\]
S'il existe une involution du couplage qui échange les deux lecteurs, on utilise \(L_{\rm sym}\). Sinon, on présente leur mesure spectrale conjointe.

\Needspace{7\baselineskip}
\subsection{Codomaine du lecteur cellulaire : caractère adimensionnel, tour et réalisation de masse}
Le résultat peut être :

\[
 \begin{cases}
 m_X\in\mathcal L_M,&\text{masse scalaire},\\
 \operatorname{Spec}(\widehat M_X),&\text{multiplet},\\
 z_X=M_X-i\Gamma_X/2,&\text{résonance},\\
 0\in\mathcal L_M,&\text{section nulle},\\
 (\mathcal C,F,R,\Delta E),&\text{secteur topologique},\\
 (L,\operatorname{Obs}_L),&\text{section virtuelle}.
 \end{cases}
\]
Cette diversité de codomaines fait partie de la loi ; elle n'est pas une exception.

\Needspace{7\baselineskip}
